%----------------------------------------------------------------------------------------------------
% START OF CHAPTER
%----------------------------------------------------------------------------------------------------
\part{Quick reference}

\chapter{Common Pitfalls}
\section{Basic slip-ups}
\begin{table}
    \centering
    \begin{tabular}{c | c | l }
        \textbf{Correct}          & \textbf{Wrong!}      &  \textit{Reason} \\
        \hline
        \itb{ho caldo}            & sono caldo      & \textit{Feelings use \iti{avere}}  \\
        \itb{ho 20 anni}          & sono 20 anni    & \textit{Age uses \iti{avere}}  \\
        \itb{non vado mai}        & non mai vado    & \textit{\iti{Non} is before verb, \iti{mai} after}  \\
        \itb{si sono alzate}      & sono alzate     & \textit{Reflexive verb needs pronoun}  \\
        \itb{siamo andati}        & abbiamo andato  & \textit{\iti{Andare} uses \iti{essere} in perfect tense}  \\
        \itb{sono le dodici}      & è dodici        & \textit{Time uses \iti{sono le}}  \\
        \itb{è l'una}             & è l'uno         & \textit{One o'clock is feminine}  \\
        \itb{il quindici marzo}   & il quindici di marzo   &  \textit{dates do not use \iti{di}, just number and month}  \\
        \itb{nella nostra scuola} & in la nostra scuola    & \textit{\iti{In} + \iti{la} = \iti{nella}}  \\
        \itb{ogni giorno}         & ogni giorni     & \textit{\iti{giorno} used with \iti{ogni} is singular}  \\
        \itb{arancione}           & arancioni       & \textit{One of the invariant color names, like \iti{blu, rosa, viola}}  \\
        \itb{simpatici}           & simpatichi      & \textit{No \textbf{h} in masculine plural, only feminine}  \\
        \itb{le notti}            & le notte        & \textit{Nouns ending \textbf{-e} take \textbf{-i} in plural}  \\
        \itb{gli amici}           & gl'amici        & \textit{\iti{Gli} never contracts}  \\
        \itb{voglio vederlo}      & voglio lo vedere       & \textit{Pronoun attaches to infinitive}  \\
        \itb{è venuto}            & ha venuto       & \textit{\iti{Venire} uses \iti{essere} in perfect tense}  \\
        \itb{è stata bellissima}  & era bellissima  & \textit{Use the perfect tense for completed periods}  \\
    \end{tabular}
\end{table}

\section{More subtle errors}
\begin{table}
    \centering
    \begin{tabular}{c | l}
        \textbf{What}        & \textbf{Comment}  \\
        \hline
        \itb{tutto/tutti}    & Contrast \iti{tutto} for every\uline{thing}, with \iti{tutti} for every\uline{one} \\
        \itb{chi/che}        & Use \iti{chi} for who, but \iti{che} for `who' (`that') as a relative clause \\
        \itb{non} + \itb{nessuno}  & Use \iti{nessuno} alone if it is the subject. If it is the object, put \iti{non} before the verb \\
        \itb{magari}         & When using \iti{magari}, `if only', follow it with the subjunctive tense \\
        \itb{città}          & The plural of \iti{città} is... \iti{città}
    \end{tabular}
\end{table}

\section{Examples}
\begin{itemize}
    \item \iti{Gli studenti che studiano l'italiano ogni giorno capiscono tutto: chi non studia mai non capisce niente!} \\
               The students (who/that) study Italian every day understand everything: (those who) never study understand nothing!

    \item \iti{In inverno i miei fratelli e le mie sorelle non escono mai: restano sempre a casa, guardano la televisione e vanno a letto presto.} \\
               In winter my brothers and sisters never go out: they always stay home, watch TV and go to bed early.

    \item \iti{Perché non mi hai chiamato? Ho aspettato tutta la sera e nessuno è venuto!} \\
               Why didn't you call me? I waited all evening and nobody called!

    \item \iti{Non hai mai visitato il Giappone? Sono andato l'anno scorso ed era il paese più bello del mondo.} \\
               Have you ever visited Japan? I went last year and it was the most beautiful country in the world.

    \item \iti{Che ora è? Sono le dodici meno un quarto. È l'una. È mezzanotte. Sono le sei e un quarto} \\
               What time is it? It's a quarter to twelve. It's one o'clock. It's midnight. It's 6:15.

    \item \iti{Nessuno mi ha chiamato. Non ho visto nessuno.} \\
               Nobody called me. I haven't seen anybody.

    \item \iti{Chi è venuto alla tua festa il sabato scorso e cosa hanno detto?} \\
               Who came to your party last Saturday and what did they say?

    \item \iti{I miei fratelli e le mie sorelle non sono mai andati in Italia. Magari potessero venire con noi.} \\
               My brothers and sisters have never been to Italy. If only they could come with me.

    \item \iti{La domenica le mie sorelle non fanno mai niente - non si alzano, non si vestono e non escono!} \\
               On Sundays my sisters never do anything - they don't get up, they don't get dressed, and they don't go out!

    \item \iti{Le studentesse si sono alzate in ritardo, non si sono vestite e non sono uscite - non fanno mai niente la domenica!} \\
               The students got up late, didn't dress and did not go out. They never do anything on Sundays!

    \item \iti{Le mie sorelle si alzano presto ogni giorno, si vestono rapidamente ed escono alle sette e un quarto. Devono essere a scuola alle sette e mezza} \\
               My sisters get up early every day, dress quickly and leave at 7:15. They need to be at school at 7:30.

    \item \iti{L'estate scorsa è stata fantastica. Siamo andati in Italia, abbiamo visitato cinque città,
               ci siamo divertiti molto e non abbiamo mai voluto tornare a casa. Magari potessimo tornare!} \\
               Last summer was wonderful. We went to Italy, visited five cities, had a lot of fun and never wanted to come home. If only we could return!

\end{itemize}
\chapter{Verb Conjugation Quick Reference}

\section{Present Tense Summary}
\begin{table}[h]
    \centering
    \begin{tabular}{r c c c c c c }
        \textbf{Verb} & \iti{io} & \iti{tu} & \iti{lui/lei} & \iti{noi} & \iti{voi} & \iti{loro} \\
        \hline
        \iti{essere} &  \iti{sono}  &  \iti{sei}  &  \iti{è}  &  \iti{siamo}  &  \iti{siete}  &  \iti{sono}  \\
        \iti{avere}  &  \iti{ho}  &  \iti{hai}  &  \iti{ha}  &  \iti{abbiamo}  &  \iti{avete}  &  \iti{hanno}  \\
        \iti{stare}  &  \iti{sto}  &  \iti{stai}  &  \iti{sta}  &  \iti{stiamo}  &  \iti{state}  &  \iti{stanno}  \\
        \iti{andare} &  \iti{vado}  &  \iti{vai}  &  \iti{va}  &  \iti{andiamo}  &  \iti{andate}  &  \iti{vanno}  \\
        \iti{fare}   &  \iti{faccio}  &  \iti{fai}  &  \iti{fa}  &  \iti{facciamo}  &  \iti{fate}  &  \iti{fanno}  \\
        \iti{venire} &  \iti{vengo}  &  \iti{vieni}  &  \iti{viene}  &  \iti{veniamo}  &  \iti{venite}  &  \iti{vengono}  \\
        \iti{uscire} &  \iti{esco}  &  \iti{esci}  &  \iti{esce}  &  \iti{usciamo}  &  \iti{uscite}  &  \iti{escono}  \\
        \iti{potere} &  \iti{posso}  &  \iti{puoi}  &  \iti{può}  &  \iti{possiamo}  &  \iti{potete}  &  \iti{possono}  \\
        \iti{volere} &  \iti{voglio}  &  \iti{vuoi}  &  \iti{vuole}  &  \iti{vogliamo}  &  \iti{volete}  &  \iti{vogliono}  \\
        \iti{dovere} &  \iti{devo}  &  \iti{devi}  &  \iti{deve}  &  \iti{dobbiamo}  &  \iti{dovete}  &  \iti{devono}  \\
        \iti{sapere} &  \iti{so}  &  \iti{sai}  &  \iti{sa}  &  \iti{sappiamo}  &  \iti{sapete}  &  \iti{sanno}  \\
    \end{tabular}
\end{table}


\section{Gerund Summary}
\begin{table}[h]
    \centering
    \begin{tabular}{r l }
      \textbf{Verb} & \textbf{Gerund}      \\
      \hline
      \iti{parlare} & \iti{parlando}  \\
      \iti{leggere} & \iti{leggendo}  \\
      \iti{dormire} & \iti{dormendo}  \\
      \iti{fare}    & \iti{facendo}   \\
      \iti{bere}    & \iti{bevendo}   \\
      \iti{dire}    & \iti{dicendo}   \\
    \end{tabular}
\end{table}


\section{Past Participle Quick Reference}
\begin{table}[h]
    \centering
    \begin{tabular}{c c l}
        \hline
        \iti{parlare}     &   \iti{parlato}          &  \iti{avere}  \\
        \iti{mangiare}    &   \iti{mangiato}         &  \iti{avere}  \\
        \iti{vedere}      &   \iti{visto}            &  \iti{avere}  \\
        \iti{leggere}     &   \iti{letto}            &  \iti{avere}  \\
        \iti{scrivere}    &   \iti{scritto}          &  \iti{avere}  \\
        \iti{fare}        &   \iti{fatto}            &  \iti{avere}  \\
        \iti{dire}        &   \iti{detto}            &  \iti{avere}  \\
        \iti{bere}        &   \iti{bevuto}           &  \iti{avere}  \\
        \iti{prendere}    &   \iti{preso}            &  \iti{avere}  \\
        \iti{perdere}     &   \iti{perso}            &  \iti{avere}  \\
        \iti{aprire}      &   \iti{aperto}           &  \iti{avere}  \\
        \iti{andare}      &   \iti{andato}           &  \iti{essere}  \\
        \iti{venire}      &   \iti{venuto}           &  \iti{essere}  \\
        \iti{uscire}      &   \iti{uscito}           &  \iti{essere}  \\
        \iti{partire}     &   \iti{partito}          &  \iti{essere}  \\
        \iti{arrivare}    &   \iti{arrivato}         &  \iti{essere}  \\
        \iti{tornare}     &   \iti{tornato}          &  \iti{essere}  \\
        \iti{nascere}     &   \iti{nato}             &  \iti{essere}  \\
        \iti{essere}      &   \iti{stato}            &  \iti{essere}  \\
        \iti{alzarsi}     &   \iti{alzato/a/i/e}     &  \iti{essere} (reflexive)  \\
        \iti{divertirsi}  &   \iti{divertito/a/i/e}  &  \iti{essere} (reflexive)  \\
    \end{tabular}
\end{table}

%----------------------------------------------------------------------------------------------------
% END OF CHAPTER
%----------------------------------------------------------------------------------------------------